% GENERATED FILE. Edit canonical lessons, then run npm run generate:books.
% canonical-source-hash: fnv1a64:5da5f6cffb8eee4d
% canonical-lessons: LA-C516-guttur, LA-C516-fauces, LA-C516-cervix, LA-C516-membrum, LA-C516-artus

\chapter{Throat, neck and limbs}
\label{ch:la-throat-neck-and-limbs}
\hlchaptermodality{516}

\begin{chapteropening}
\textbf{By the end of this chapter:} \emph{I can name the throat, the neck and the limbs.}

Complete the last lesson of chapter 516: I can name the throat, the neck and the limbs.
\end{chapteropening}

\section[guttur, gutturis]{guttur, gutturis — throat, gullet}

\label{lesson:LA-C516-guttur}



\begin{quote}
Before the new one: say the Latin for eyebrow, then the Latin for eyelid.
\end{quote}

\subsection*{You'll want to know: guttur, gutturis}
\textbf{guttur, gutturis} — \textquotedblleft{}throat, gullet\textquotedblright{}.

\begin{quote}
Guttur mihi dolet.
\end{quote}

\begin{quote}
\emph{My throat hurts.}
\end{quote}

Neuter, third declension. English \emph{guttural} — of the throat.

\begin{cousinweb}
\textbf{guttur} is inherited; EtymDB traces it to a Proto-Indo-European root meaning \textquotedblleft{}a swelling, rounding\textquotedblright{}. English \emph{guttural} comes from it.
\end{cousinweb}

\begin{grammarlens}[title={What you've built}]
One more word for telling, explaining and describing at B1.
\end{grammarlens}

\subsection*{Your turn}
\begin{itemize}
\raggedright
  \item \emph{Say it:} \emph{guttur, gutturis}
  \item \emph{Say it:} \emph{guttur, gutturis}, once more
  \item \emph{Say it:} a sentence of your own with \emph{guttur, gutturis}
  \item \emph{Recall:} say the Latin for eyebrow, then the Latin for eyelid, then say \emph{guttur, gutturis} again
\end{itemize}

\begin{tcolorbox}[breakable,colback=teal!4,colframe=teal!35!black,title={Before you move on}]
What does guttur, gutturis mean? (\textquotedblleft{}Throat, gullet\textquotedblright{}.) Say it once more.
\end{tcolorbox}
\section[faucēs, faucium]{faucēs, faucium — jaws, the back of the mouth; a narrow pass}

\label{lesson:LA-C516-fauces}



\begin{quote}
Before the new one: say the Latin for eyelid, then the Latin for throat.
\end{quote}

\subsection*{You'll want to know: faucēs, faucium}
\textbf{faucēs, faucium} — \textquotedblleft{}jaws, the back of the mouth; a narrow pass\textquotedblright{}.

\begin{quote}
Leō faucēs aperit.
\end{quote}

\begin{quote}
\emph{The lion opens its jaws.}
\end{quote}

Feminine, third declension, plural. The back of the mouth; also a narrow entrance to a valley or a harbour.

\begin{cousinweb}
\textbf{faucēs} is a native Latin noun; Lewis \& Short compares a Sanskrit word for \textquotedblleft{}opening\textquotedblright{}, but the origin is not settled.
\end{cousinweb}

\begin{grammarlens}[title={What you've built}]
One more word for telling, explaining and describing at B1.
\end{grammarlens}

\subsection*{Your turn}
\begin{itemize}
\raggedright
  \item \emph{Say it:} \emph{faucēs, faucium}
  \item \emph{Say it:} \emph{faucēs, faucium}, once more
  \item \emph{Say it:} a sentence of your own with \emph{faucēs, faucium}
  \item \emph{Recall:} say the Latin for eyelid, then the Latin for throat, then say \emph{faucēs, faucium} again
\end{itemize}

\begin{tcolorbox}[breakable,colback=teal!4,colframe=teal!35!black,title={Before you move on}]
What does faucēs, faucium mean? (\textquotedblleft{}Jaws, the back of the mouth; a narrow pass\textquotedblright{}.) Say it once more.
\end{tcolorbox}
\section[cervīx, cervīcis]{cervīx, cervīcis — neck, nape}

\label{lesson:LA-C516-cervix}



\begin{quote}
Before the new one: say the Latin for throat, then the Latin for jaws.
\end{quote}

\subsection*{You'll want to know: cervīx, cervīcis}
\textbf{cervīx, cervīcis} — \textquotedblleft{}neck, nape\textquotedblright{}.

\begin{quote}
Servus sarcinam in cervīcibus portat.
\end{quote}

\begin{quote}
\emph{The slave carries the pack on his neck.}
\end{quote}

Feminine, third declension; often plural \emph{cervīcēs}. English \emph{cervical}.

\begin{cousinweb}
\textbf{cervīx} is inherited; EtymDB traces it to a Proto-Indo-European root meaning \textquotedblleft{}horn, head\textquotedblright{}. English \emph{cervical} comes from it.
\end{cousinweb}

\begin{grammarlens}[title={What you've built}]
One more word for telling, explaining and describing at B1.
\end{grammarlens}

\subsection*{Your turn}
\begin{itemize}
\raggedright
  \item \emph{Say it:} \emph{cervīx, cervīcis}
  \item \emph{Say it:} \emph{cervīx, cervīcis}, once more
  \item \emph{Say it:} a sentence of your own with \emph{cervīx, cervīcis}
  \item \emph{Recall:} say the Latin for throat, then the Latin for jaws, then say \emph{cervīx, cervīcis} again
\end{itemize}

\begin{tcolorbox}[breakable,colback=teal!4,colframe=teal!35!black,title={Before you move on}]
What does cervīx, cervīcis mean? (\textquotedblleft{}Neck, nape\textquotedblright{}.) Say it once more.
\end{tcolorbox}
\section[membrum, membrī]{membrum, membrī — a limb; a member}

\label{lesson:LA-C516-membrum}



\begin{quote}
Before the new one: say the Latin for jaws, then the Latin for neck.
\end{quote}

\subsection*{You'll want to know: membrum, membrī}
\textbf{membrum, membrī} — \textquotedblleft{}a limb; a member\textquotedblright{}.

\begin{quote}
Post longum iter membra mea dolent.
\end{quote}

\begin{quote}
\emph{After the long journey my limbs ache.}
\end{quote}

Neuter, second declension. Also a member of a group. English \emph{member}.

\begin{cousinweb}
\textbf{membrum} is inherited through Proto-Italic from a Proto-Indo-European word for \textquotedblleft{}flesh\textquotedblright{} (EtymDB). English \emph{member} comes from it.
\end{cousinweb}

\begin{grammarlens}[title={What you've built}]
One more word for telling, explaining and describing at B1.
\end{grammarlens}

\subsection*{Your turn}
\begin{itemize}
\raggedright
  \item \emph{Say it:} \emph{membrum, membrī}
  \item \emph{Say it:} \emph{membrum, membrī}, once more
  \item \emph{Say it:} a sentence of your own with \emph{membrum, membrī}
  \item \emph{Recall:} say the Latin for jaws, then the Latin for neck, then say \emph{membrum, membrī} again
\end{itemize}

\begin{tcolorbox}[breakable,colback=teal!4,colframe=teal!35!black,title={Before you move on}]
What does membrum, membrī mean? (\textquotedblleft{}A limb; a member\textquotedblright{}.) Say it once more.
\end{tcolorbox}
\section[artus, artuum]{artus, artuum — joints, limbs}

\label{lesson:LA-C516-artus}



\begin{quote}
Before the new one: say the Latin for neck, then the Latin for a limb.
\end{quote}

\subsection*{You'll want to know: artus, artuum}
\textbf{artus, artuum} — \textquotedblleft{}joints, limbs\textquotedblright{}.

\begin{quote}
Frīgus artūs senis rigidōs facit.
\end{quote}

\begin{quote}
\emph{The cold makes the old man's joints stiff.}
\end{quote}

Masculine, fourth declension, nearly always plural; dative and ablative \emph{artubus}. English \emph{article}, a little joint, comes from its diminutive.

\begin{cousinweb}
\textbf{artus} belongs with \emph{arma} and a Proto-Indo-European root meaning \textquotedblleft{}to fit together\textquotedblright{} (Lewis \& Short; EtymDB). English \emph{article} comes from its diminutive \emph{articulus} \textquotedblleft{}a little joint\textquotedblright{}.
\end{cousinweb}

\begin{grammarlens}[title={What you've built}]
One more word for telling, explaining and describing at B1.
\end{grammarlens}

\subsection*{Your turn}
\begin{itemize}
\raggedright
  \item \emph{Say it:} \emph{artus, artuum}
  \item \emph{Say it:} \emph{artus, artuum}, once more
  \item \emph{Say it:} a sentence of your own with \emph{artus, artuum}
  \item \emph{Recall:} say the Latin for neck, then the Latin for a limb, then say \emph{artus, artuum} again
\end{itemize}

\begin{tcolorbox}[breakable,colback=teal!4,colframe=teal!35!black,title={Before you move on}]
What does artus, artuum mean? (\textquotedblleft{}Joints, limbs\textquotedblright{}.) Say it once more.
\end{tcolorbox}
